% ===========================================================================
\chapter{Especificação dos Casos de Uso de Gestão Pedagógica de Turmas}
\label{cap_gestao_turmas}
% ===========================================================================

A gestão acadêmica no SIGEA-GTT é orientada pelo modelo de docência titular, no qual o professor é a autoridade pedagógica responsável pela criação das turmas, alocação dos estudantes e publicação das atividades de auditoria clínica com prontuários simulados. A Figura \ref{fig_uc_modulo_turmas} exibe o diagrama de casos de uso do subsistema em notação UML com renderização vetorial TikZ:

\begin{figure}[!ht]
\centering
\caption{Diagrama de Casos de Uso do Subsistema de Gestão Acadêmica}
\label{fig_uc_modulo_turmas}
\resizebox{0.90\textwidth}{!}{%
\begin{tikzpicture}[
    actor/.style={
        draw=clinical-900,
        fill=clinical-100!40,
        thick,
        rounded corners=2pt,
        align=center,
        font=\bfseries\scriptsize,
        minimum width=2.4cm,
        minimum height=1.2cm
    },
    usecase/.style={
        ellipse,
        draw=clinical-700,
        fill=white,
        thick,
        align=center,
        font=\scriptsize,
        minimum width=3.4cm,
        minimum height=1.1cm,
        inner sep=2pt
    },
    link/.style={
        draw=clinical-900!80,
        thick
    }
]
    % Fronteira
    \draw[rounded corners=8pt, fill=clinical-100!15, draw=clinical-700, line width=1.2pt] 
        (0.5, -5.5) rectangle (12.5, 3.2);
    \node[anchor=north west, font=\bfseries\small\color{clinical-700}] at (0.8, 2.9) 
        {Subsistema: Gestão Pedagógica de Turmas};

    % Atores
    \node[actor] (prof) at (-2.2, 0.5) {Docente Titular\\\texttt{\scriptsize (PROFESSOR)}};
    \node[actor] (admin) at (-2.2, -3.5) {Administrador\\\texttt{\scriptsize (ADMIN)}};

    % Casos de Uso
    \node[usecase] (uc07) at (3.5, 1.8) {UC07: Criar e Configurar\\Turma Acadêmica};
    \node[usecase] (uc08) at (8.5, 1.8) {UC08: Matricular e Desmatricular\\Discentes na Turma};
    \node[usecase] (uc09) at (3.5, -1.0) {UC09: Cadastrar Atividade\\com Prontuário Simulado};
    \node[usecase] (uc10) at (8.5, -1.0) {UC10: Encerrar Período\\Letivo da Turma};

    % Conexões
    \draw[link] (prof) -- (uc07);
    \draw[link] (prof) -- (uc08);
    \draw[link] (prof) -- (uc09);
    \draw[link] (prof) -- (uc10);

    \draw[link] (admin) -- (uc07);
    \draw[link] (admin) -- (uc10);

\end{tikzpicture}
}
\fonte{Elaborado pelo autor (2026).}
\end{figure}

% ===========================================================================
\section{Especificação Detalhada dos Casos de Uso}
\label{sec_especificacao_uc_turmas}
% ===========================================================================

% ---------------------------------------------------------------------------
% UC07
% ---------------------------------------------------------------------------
\subsection{UC07: Criar e Configurar Turma Acadêmica}
\label{uc07_especificacao}

\begin{table}[!ht]
\centering
\caption{Especificação Detalhada do Caso de Uso UC07}
\label{tab_uc07}
\footnotesize
\begin{tabularx}{\textwidth}{p{3.5cm} X}
\toprule
\textbf{Propriedade} & \textbf{Especificação Técnica e Comportamental} \\
\midrule
\textbf{Identificador e Nome} & \texttt{UC07: Criar e Configurar Turma Acadêmica} \\
\textbf{Atores} & \texttt{PROFESSOR} (Ator Primário), \texttt{ADMIN} (Ator Secundário). \\
\textbf{Nível de Meta} & Nível Usuário (\textit{Sea Level}). \\
\textbf{Pré-condições} & Usuário autenticado com perfil docente (\texttt{PROFESSOR}) ou administrativo. \\
\textbf{Gatilho} & O professor acessa o painel "Minhas Turmas" e clica em "Nova Turma". \\
\textbf{Fluxo Principal} & 
1. O sistema exibe o formulário de cadastro de turma.\\
& 2. O professor preenche o nome da disciplina (ex.: \textit{"Enfermagem na Saúde do Adulto I"}), o código da turma (ex.: \textit{"T01"}) e o período letivo semestral (ex.: \textit{"2025.2"}).\\
& 3. O sistema associa automaticamente o ID do professor logado ao campo \texttt{professor\_id}.\\
& 4. O sistema valida se o formato de exibição composto atende ao padrão \texttt{"Disciplina - Turma - Período"}.\\
& 5. O sistema verifica a unicidade da combinação \texttt{(subject\_name, class\_code, academic\_period)}.\\
& 6. O sistema persiste o registro na tabela \texttt{academic\_classes} com \texttt{is\_closed = false}.\\
& 7. O sistema exibe notificação de confirmação e adiciona a turma à lista do professor. \\
\textbf{Fluxos de Exceção} & 
\textbf{[5e] Turma Duplicada}: Se já existir uma turma com o mesmo nome, código e período letivo, o sistema bloqueia o cadastro e emite alerta de duplicidade.\\
& \textbf{[3e] Ausência de Professor Titular}: A tentativa de submissão sem vínculo docente é barrada pela restrição \texttt{NOT NULL} do modelo de dados. \\
\textbf{Pós-condições} & Turma acadêmica criada e disponível para alocação de discentes e atividades. \\
\textbf{Regras de Negócio} & \textbf{RN06} (Padrão de Nomeclatura de Turmas) e \textbf{RN07} (Titularidade Docente Obrigatória). \\
\textbf{Requisitos Associados} & \texttt{RF10}, \texttt{RNF01}, \texttt{RNF03}. \\
\bottomrule
\end{tabularx}
\fonte{Elaborado pelo autor (2026).}
\end{table}

% ---------------------------------------------------------------------------
% UC08
% ---------------------------------------------------------------------------
\subsection{UC08: Matricular e Desmatricular Discentes na Turma}
\label{uc08_especificacao}

\begin{table}[!ht]
\centering
\caption{Especificação Detalhada do Caso de Uso UC08}
\label{tab_uc08}
\footnotesize
\begin{tabularx}{\textwidth}{p{3.5cm} X}
\toprule
\textbf{Propriedade} & \textbf{Especificação Técnica e Comportamental} \\
\midrule
\textbf{Identificador e Nome} & \texttt{UC08: Matricular e Desmatricular Discentes na Turma} \\
\textbf{Atores} & \texttt{PROFESSOR} (Ator Primário). \\
\textbf{Nível de Meta} & Nível Usuário (\textit{Sea Level}). \\
\textbf{Pré-condições} & Turma criada e aberta (\texttt{is\_closed = false}). Discentes cadastrados no sistema com perfil \texttt{STUDENT}. \\
\textbf{Gatilho} & O professor abre os detalhes da turma e acessa a aba "Discentes Matriculados". \\
\textbf{Fluxo Principal} & 
1. O sistema apresenta a lista de estudantes já matriculados na turma.\\
& 2. O professor aciona o botão "Adicionar Estudantes".\\
& 3. O sistema abre um modal com a lista de todos os usuários ativos com perfil \texttt{STUDENT}, exibindo nome, e-mail institucional e número de matrícula.\\
& 4. O professor seleciona um ou múltiplos alunos por meio de caixas de seleção.\\
& 5. O professor confirma a inclusão.\\
& 6. O sistema valida que nenhum dos alunos selecionados possui matrícula prévia na mesma turma.\\
& 7. O sistema insere os registros na tabela associativa \texttt{class\_students} com carimbo temporal \texttt{enrolled\_at}.\\
& 8. O sistema atualiza o contador de alunos da turma e emite feedback \textit{Toast}. \\
\textbf{Fluxos Alternativos} & 
\textbf{[1a] Desmatrícula de Discente}: O professor localiza um estudante na lista da turma e clica no ícone de exclusão. O sistema abre modal de confirmação. Após a confirmação, o registro na tabela \texttt{class\_students} é removido. \\
\textbf{Fluxos de Exceção} & 
\textbf{[1e] Turma Fechada}: Se \texttt{is\_closed = true}, as opções de inclusão e remoção de alunos permanecem desabilitadas na interface.\\
& \textbf{[6e] Matrícula em Duplicidade}: A chave primária composta impede a reinserção de um mesmo estudante na turma. \\
\textbf{Pós-condições} & Relação N:N atualizada entre a turma e o corpo discente. \\
\textbf{Regras de Negócio} & Integridade de matrícula e segregação por perfil discente. \\
\textbf{Requisitos Associados} & \texttt{RF11}, \texttt{RNF03}. \\
\bottomrule
\end{tabularx}
\fonte{Elaborado pelo autor (2026).}
\end{table}

% ---------------------------------------------------------------------------
% UC09
% ---------------------------------------------------------------------------
\subsection{UC09: Cadastrar Atividade Prática com Prontuário Simulado}
\label{uc09_especificacao}

\begin{table}[!ht]
\centering
\caption{Especificação Detalhada do Caso de Uso UC09}
\label{tab_uc09}
\footnotesize
\begin{tabularx}{\textwidth}{p{3.5cm} X}
\toprule
\textbf{Propriedade} & \textbf{Especificação Técnica e Comportamental} \\
\midrule
\textbf{Identificador e Nome} & \texttt{UC09: Cadastrar Atividade Prática com Prontuário Simulado} \\
\textbf{Atores} & \texttt{PROFESSOR} (Ator Primário). \\
\textbf{Nível de Meta} & Nível Usuário (\textit{Sea Level}). \\
\textbf{Pré-condições} & O docente titular deve possuir turma aberta e ativa. \\
\textbf{Gatilho} & O professor clica em "Nova Atividade" no menu de sua turma. \\
\textbf{Fluxo Principal} & 
1. O sistema apresenta o formulário com título, descrição instrucional, data/hora limite de entrega (\texttt{deadline}) e editor do caso clínico simulado.\\
& 2. O professor preenche o título (ex.: \textit{"Auditoria GTT -- Caso Cirúrgico Paciente Idoso"}).\\
& 3. O professor define a data limite de entrega posterior ao horário atual.\\
& 4. O professor insere os dados do caso clínico estruturado (dados sociodemográficos fictícios, queixa principal, evolução médica, prescrições ativas/suspensas, anotações de enfermagem e exames laboratoriais com valores de referência).\\
& 5. O sistema valida o esquema JSON do prontuário simulado (\texttt{clinical\_case\_data}).\\
& 6. O sistema persiste a atividade na tabela \texttt{activities} vinculada à turma.\\
& 7. A atividade torna-se visível no painel dos discentes matriculados naquela turma. \\
\textbf{Fluxos de Exceção} & 
\textbf{[3e] Prazo Limite Inválido}: Se a data limite for anterior à data e hora do cadastro, o sistema rejeita o formulário com mensagem de erro.\\
& \textbf{[5e] Esquema JSON Clínico Inválido}: Se faltarem nós essenciais no JSONB (ex.: ausência do nó \texttt{patient} ou \texttt{prescriptions}), o validador do backend rejeita a submissão. \\
\textbf{Pós-condições} & Atividade cadastrada e pronta para ser executada pelos alunos da turma. \\
\textbf{Regras de Negócio} & Isolamento de dados clínicos simulados e prazo cronológico válido. \\
\textbf{Requisitos Associados} & \texttt{RF12}, \texttt{RF14}, \texttt{RNF02}. \\
\bottomrule
\end{tabularx}
\fonte{Elaborado pelo autor (2026).}
\end{table}

% ---------------------------------------------------------------------------
% UC10
% ---------------------------------------------------------------------------
\subsection{UC10: Encerrar Período Letivo da Turma}
\label{uc10_especificacao}

\begin{table}[!ht]
\centering
\caption{Especificação Detalhada do Caso de Uso UC10}
\label{tab_uc10}
\footnotesize
\begin{tabularx}{\textwidth}{p{3.5cm} X}
\toprule
\textbf{Propriedade} & \textbf{Especificação Técnica e Comportamental} \\
\midrule
\textbf{Identificador e Nome} & \texttt{UC10: Encerrar Período Letivo da Turma} \\
\textbf{Atores} & \texttt{PROFESSOR} (Ator Primário), \texttt{ADMIN} (Ator Secundário). \\
\textbf{Nível de Meta} & Nível Usuário (\textit{Sea Level}). \\
\textbf{Pré-condições} & A turma está ativa (\texttt{is\_closed = false}). \\
\textbf{Gatilho} & O docente clica no botão "Encerrar Turma" nas opções de gestão. \\
\textbf{Fluxo Principal} & 
1. O sistema aciona rotina de verificação de consistência acadêmica.\\
& 2. O sistema checa se todas as submissões enviadas pelos discentes nas atividades daquela turma já receberam avaliação e nota docente (\texttt{grade IS NOT NULL}).\\
& 3. Não existindo pendências de correção, o sistema abre modal de confirmação avisando sobre o congelamento irrevogável da turma.\\
& 4. O professor confirma a operação.\\
& 5. O sistema atualiza o atributo \texttt{is\_closed} para \texttt{true}.\\
& 6. O sistema bloqueia a criação de novas atividades, matrículas e submissões tardias, mantendo os dados disponíveis apenas para consulta histórica. \\
\textbf{Fluxos de Exceção} & 
\textbf{[2e] Pendência de Avaliação}: Caso existam submissões de auditoria entregues por alunos que ainda não foram avaliadas pelo professor (\texttt{grade IS NULL}), o sistema aborta a operação, exibe aviso de pendência pedagógica e lista as atividades e discentes não avaliados. \\
\textbf{Pós-condições} & Turma encerrada e congelada contra novas modificações. \\
\textbf{Regras de Negócio} & \textbf{RN08} (Bloqueio de Encerramento com Pendências de Correção). \\
\textbf{Requisitos Associados} & \texttt{RF13}, \texttt{RNF03}. \\
\bottomrule
\end{tabularx}
\fonte{Elaborado pelo autor (2026).}
\end{table}
